% 395_Messkette_Kanal_En_ch.tex
% Johann Pascher, 24 Oct 2026
% The measurement chain as a channel

\providecommand{\BM}{\textbf{[B]}}
\providecommand{\KM}{\textbf{[K]}}
\providecommand{\SM}{\textbf{[S]}}
\providecommand{\XM}{\textbf{[X]}}
\providecommand{\QM}{\textbf{[Q]}}

\chapter*{The Measurement Chain as a Channel}
\addcontentsline{toc}{chapter}{The Measurement Chain as a Channel}

\begin{abstract}
Communications engineering no longer fights echoes; it measures and uses them: known pilot symbols
reveal the channel, a suitable basis makes it diagonal, regeneration prevents noise from adding up,
and the Shannon limit states how much information a noisy channel can deliver at most. This document
transfers the four methods to the question of whether recursions are part of the measured values
(Doc.~394). Results: in the ratio sector the channel is flat; with the Koide form as pilot, the lepton
ratios deviate by less than $1\,\xi$. A mode-dependent channel factor is thereby limited to at most
about $1\,\%$ of the correction step $100\,\xi$. In the SI sector there is exactly one pilot,
$\alpha$, and there cannot be more in a ratio-based theory: the only relation with an SI unit is
$\alpha_{\text{SI}}=\xi\,(E_0/1\,\text{MeV})^2$, and $K$ and the MeV anchoring are the same degree of
freedom. The muon lifetime tests not $K$ but the dimensionless ratio $v/m_\mu$; the bound on the mode
factor needs only the derived angle $2/9$. For quantum computers it follows that a common gate factor is absorbed
by the calibration and is in principle not measurable; only a mode-dependent part of about $10^{-4}$
would be visible, far below today's decoherence. Discrete structures with small
denominators survive the measurement chain; the SI correction of $1.33\,\%$ leaves denominators
unambiguous only up to about~$11$. As a yardstick against number mysticism, a hit counts only if the
measurement delivers more information than the formula contains; the Koide parameters $2/9$ and
$\sqrt2$ pass this clearly, and so does $\alpha^{-1}=3700/27$ when it is calculated according to its
constituents $7500\cdot(74/75)/54$ rather than as a freely searched fraction.
\end{abstract}

% ============================================================
\section*{Status markers}
% ============================================================
\begin{center}
\begin{tabular}{@{}lp{11cm}@{}}
\textbf{[B]} & Algebraically proved (rearrangement without additional assumption). \\[2pt]
\textbf{[K]} & Numerically confirmed (check script). \\[2pt]
\textbf{[S]} & Assumption, open, or classification without a calculation of its own. \\[2pt]
\textbf{[X]} & Excluded by the calculation. \\[2pt]
\textbf{[Q]} & Taken from the literature, with source. \\
\end{tabular}
\end{center}
Measured values: $m_e$, $m_\mu$ and $\alpha$ from CODATA 2022 \QM, $m_\tau=1776.93\pm0.09$\,MeV from
PDG (pdgLive, 2026 edition) \QM.

% ============================================================
\section{Occasion}
% ============================================================

The basic form of FFGFT calculates with ratios; a fractal correction appears only in the translation
into SI units. Behind this stands the reading that the measurer stands inside the scale recursion, so
that the recursion is part of every measured value. Doc.~394 showed that this reading has so far been
tested on exactly one quantity, $\alpha$, and that the other residuals have a different origin. What
remained open was how to test the question systematically.

Communications engineering has faced the same situation for decades: a signal runs through a channel
that alters it, and the receiver stands at the end of this channel, not outside it. How it deals with
this yields four transferable methods. Two are directly usable as methods, channel estimation with a
pilot and the Shannon balance; two are structural hints, the eigenbasis and regeneration.

% ============================================================
\section{How communications engineering deals with echoes}
% ============================================================

Besides the direct signal, a radio channel delivers delayed copies, the echoes. They used to be
regarded as interference; today they are measured and used as additional signal energy \QM. The
\emph{RAKE receiver} (CDMA, UMTS) picks up the delayed copies individually with several
\enquote{fingers}, aligns their delays and adds them with weights \cite{PriceGreen1958}. \emph{OFDM}
(WLAN, LTE, 5G, DVB-T) spreads the signal over many slow subcarriers; a guard interval longer than the
echo delays turns the convolution with the echoes into a simple phase and amplitude rotation per
subcarrier \cite{WeinsteinEbert1971}. \emph{Maximum ratio combining} and MIMO exploit several
propagation paths as parallel channels \cite{Brennan1959}.

What makes this possible is \emph{channel estimation}: the transmitter sends along pilot symbols whose
value the receiver knows exactly; from what arrives, it determines the channel and removes it. A second
principle is \emph{regeneration}: every repeater decides anew between 0 and~1, so noise does not
accumulate along the line. What none of the methods circumvents is the \emph{Shannon limit}
$C=B\log_2(1+S/N)$ \cite{Shannon1948}. Using echoes increases the effective ratio $S/N$, and modern
error correction comes within fractions of a decibel of the limit \cite{Berrou1993,Arikan2009}, but
never beyond it. The echoes therefore do not disappear; they are measured, modelled as part of the
system and exploited, and the bound remains.

% ============================================================
\section{Channel estimation with a pilot}
% ============================================================

\subsection{Method and condition}

Transferred, this means: the measurement chain is the channel $H$ that turns the value of the basic
form $y_{\text{th}}$ into the measured value $y_{\text{meas}}=H\,y_{\text{th}}$. A pilot is a quantity
whose value is fixed algebraically and exactly in the basic form. $H$ follows from it. With this $H$
one calculates the other quantities back. If one common $H$ explains several residuals, there is a
uniform recursion; if each quantity needs its own $H$, there is none, and one then knows it for
certain. The condition for this is strict: the pilot must be exact independently. If it contains a
model error itself, the channel estimate absorbs this error and the calculation becomes circular.

\subsection{Ratio sector: the channel is flat}

The pilot in the ratio sector is the Koide form with $\theta=2/9$ and $d/c=\sqrt2$ (Docs.~353, 367),
\begin{equation}
  \sqrt{m_k}\;\propto\;1+\sqrt2\,\cos\!\Bigl(\tfrac29+\tfrac{2\pi k}{3}\Bigr).
\end{equation}
It fixes the ratios $m_\mu/m_e$ and $m_\tau/m_\mu$ without any measured value. The channel factor is
then
\begin{equation}
  H_{\mu e}-1=-0.074\,\xi,\qquad H_{\tau\mu}-1=-0.16\pm0.38\,\xi
\end{equation}
\KM, for the muon-electron ratio with a measurement uncertainty of only $2\times10^{-8}$. Both factors
lie below $1\,\xi$, more than a hundred times below the SI correction step $100\,\xi$. In the ratio
sector the channel is therefore flat: a recursion of the measurement chain is not visible there. This
is to be expected, since a common factor cancels in ratios (A040, Doc.~070).

\subsection{Pilot with a model error: the ladder}

If the lepton ladder $m_i\propto r_i\,\xi^{p_i}$ is taken as the pilot instead of the Koide form, it
lies $39.10\,\xi$ and $77.56\,\xi$ above the measured ratios \KM; each generation then requires its
own channel factor. This, however, is not a channel. Ladder and Koide form differ without any measured
value by $39.02\,\xi$ and $77.40\,\xi$ (Doc.~394) \KM: the apparent channel is the model error of the
pilot. This is exactly what the condition in Section~3.1 warns against. In the language of channel
estimation, the falling residuals $\varepsilon_i$ from Doc.~352 are the result of an inexact pilot.

\subsection{SI sector: one pilot, no test value}

In the SI sector, $\alpha$ is the only pilot without a fitted factor. It estimates the channel very
precisely: $K=\xi\,m_em_\mu\,\alpha^{-1}=0.986501$, that is $n=101.25$ in $K=1-n\,\xi$, with an
uncertainty below $3\times10^{-8}$ \KM; if the energy unit is fixed via $m_em_\mu=54$, it is
$n=100.056$ (Doc.~394) \KM. Whether the same $K$ also carries other SI quantities cannot be tested,
because a second pilot is missing: $E_P$ and $G$ depend on the fitted factor~10 in $v/E_P$
(Docs.~149, 383), and $v$ is a convention of the Standard Model (R112). Channel estimation is thus
possible, its test is not. Section~7 shows that a second SI pilot cannot exist in FFGFT at all.

% ============================================================
\section{Eigenbasis: one channel factor per mode}
% ============================================================

OFDM switches to a basis in which the channel becomes diagonal; the convolution with echoes turns into a
simple factor per subcarrier. FFGFT already has such a basis: the modes on $T^4$ with their
mode-bound time scale $\tilde T=1/m$. A measurement recursion would there be a factor $\kappa_i$ of
its own per mode, and this would lead one to expect generation-dependent corrections \SM.

The pilot from Section~3.2 limits this possibility. A difference between the channel factors of two
modes appears directly in the ratio of their masses. If the Koide form holds exactly, it follows that
\begin{equation}
  |\kappa_\mu-\kappa_e|\le0.074\,\xi,\qquad |\kappa_\tau-\kappa_\mu|\le0.92\,\xi\ (2\sigma)
\end{equation}
\KM. A mode-dependent part of the measurement recursion is thereby limited to at most about $1\,\%$ of
the correction step $100\,\xi$. The SI correction must be practically common to the three charged
leptons; this supports the uniform rule. The generation-dependent residuals of the ladder do not need
this basis; they are the distance from the Koide form (Section~3.3). The bound holds under the
assumption that the Koide form is exact; if it is not, the calculation only limits the sum of the
Koide residual and the mode factor \SM. Section~7 shows that the main part of the bound already holds
with the derived angle $2/9$ alone.

The Koide form itself can also be fitted freely to the measured values. With both parameters free,
one obtains $\theta=0.2222248\pm0.0000063$ and $d/c=1.4142089\pm0.000011$, that is $2/9$ and $\sqrt2$
each within one standard deviation, the uncertainty coming solely from the tau mass \KM.

% ============================================================
\section{Regeneration: the discrete survives the measurement chain}
% ============================================================

Digital transmission beats noise because a new decision is made at every node. This agrees with
Doc.~343: the discreteness of FFGFT is topological; it lies in winding numbers and integer structures,
not in the dynamics. Transferred, this means: exact rational ratios and integers can be recovered from
a noisy measured value as long as the tolerance is smaller than the distance to the nearest competitor.
The limit for this is the Farey height from Doc.~343,
\begin{equation}
  Q_c=\frac{\pi}{\sqrt{6\,\varepsilon}} ,
\end{equation}
up to which fractions remain unambiguous at relative tolerance $\varepsilon$:

\begin{center}
\begin{tabular}{@{}lrr@{}}
\toprule
\textbf{Tolerance} & \textbf{$\varepsilon$} & \textbf{$Q_c$}\\
\midrule
SI correction $100\,\xi$ & $1.33\,\%$ & $11.1$\\
$m_\tau/m_\mu$ & $5.1\times10^{-5}$ & $180$\\
$m_\mu/m_e$ & $2.2\times10^{-8}$ & $8700$\\
$\alpha$ & $1.5\times10^{-10}$ & $1.0\times10^{5}$\\
\bottomrule
\end{tabular}
\end{center}
\KM\ A clear rule for the presentation follows from this. If a ratio is passed through the SI
correction of $1.33\,\%$, only denominators up to about~$11$ remain unambiguous; the denominators $9$,
$5$ and $8$ of the basic form survive this, the $75$ of the rotation number does not. In the basic form
itself, without correction, the ratios are by contrast measured so precisely that large denominators
could also be recovered. Continuing corrections accumulate errors; exact ratios are preserved. This
supports the guideline of building the theory from the ratio-based basic form and applying corrections
only at the end \SM.

% ============================================================
\section{Shannon balance against number mysticism}
% ============================================================

A measurement with relative precision $\delta$ delivers, about a constant whose value is known
beforehand only to within a decade, roughly
\begin{equation}
  I=\log_2\frac{\ln 10}{2\,\delta}\ \text{bits}.
\end{equation}
A formula itself contains information; Doc.~393 measures it as complexity $C$ in bits, growing with
numerator, denominator, root and exponent. If the space of possible formulas contains more bits than the
measurement, some combination inevitably hits. A formal criterion follows from this: a hit counts only
if the surplus $I-C$ is clearly positive. The deviation from the measured value is taken as $\delta$,
but at most the measurement uncertainty.

\begin{center}
\begin{tabular}{@{}lrrr@{}}
\toprule
\textbf{Hit} & \textbf{$I$} & \textbf{$C$} & \textbf{$I-C$}\\
\midrule
$\alpha^{-1}=7400/54$ ($=3700/27$) & $17.2$ & $5.8$ & $+11.5$\\
\quad as a freely searched fraction $3700/27$ & $17.2$ & $16.6$ & $+0.6$\\
$v/E_P=\xi^4/(5\pi)$ & $9.0$ & $6.3$ & $+2.6$\\
$M_W/M_Z=\sqrt{7/9}$ & $10.8$ & $5.4$ & $+5.4$\\
$m_em_\mu$ close to the integer $54$ & $5.8$ & $0$ & $+5.8$\\
Koide $\theta=2/9$ & $15.3$ & $4.2$ & $+11.1$\\
Koide $d/c=\sqrt2$ & $17.2$ & $2.0$ & $+15.2$\\
\bottomrule
\end{tabular}
\end{center}
\KM\ For $3700/27$ the decisive point is how the formula is counted. As a freely searched fraction it
contains almost as much information as the deviation of $7.6$\,ppm provides and does not count; this
is how the null test from Doc.~393 ($p\approx0.47$) calculates. In the corpus, however, $3700/27$ is
only the reduction of $7500\cdot(74/75)/54$: $7500=1/\xi$ and $74/75$ come from the basic form, the
only new element is the integer $54$. Counted this way, more than $11$\,bits of surplus remain;
estimated directly, $7400/\alpha^{-1}$ lies this close to an integer with a probability of about
$8\times10^{-4}$ (Doc.~394) \KM. The precondition is that $74/75$ was fixed independently of $\alpha$
(Docs.~133, 300). The
closeness of the mass product to $54$ delivers just under $6$\,bits, corresponding to a chance
probability of about $1.7\,\%$ (Doc.~394). Far ahead are the two Koide parameters with $11$ and
$15$\,bits of surplus; they are the strongest numerical structure of the lepton sector. It should be
noted that $d/c=\sqrt2$ is equivalent to the Koide relation $Q=2/3$, known since 1982 \cite{Koide1982};
the surplus demonstrates the structure of the data, not yet its derivation from FFGFT. That
derivation is given for $\theta=2/9$ in Doc.~A110 and for the amplitude in Docs.~353 and 367 (R150).

The criterion is rough. It does not count how many formulas of equal complexity exist, and it depends
on the chosen bit measure \SM. For individual hits it therefore does not replace the null test from
Doc.~393, but it shows at a glance where a closer test is worthwhile and where it is not.

% ============================================================
\section{What can be resolved}
% ============================================================

\subsection{There is no second pilot in the SI sector}

The only relation in the corpus in which an SI unit appears explicitly is
$\alpha_{\text{SI}}=\xi\,(E_0/1\,\text{MeV})^2$ (Docs.~010, 011, 014, 137, 313, 365, 384, 386). All other
statements of FFGFT are pure numbers, ratios of two quantities of the same dimension; a common factor
cancels in them (A040) \BM. With $\hbar=c=1$ a single dimension remains, energy, and hence a single unit.
The SI sector therefore has exactly one degree of freedom: the factor that binds the basic form to the
MeV. $K$ and the MeV anchoring are the same number, and the one measurement $\alpha$ fixes it. A second
pilot that could test this number cannot exist in a ratio-based theory \BM. The assignment from
Doc.~386, according to which the correction sits on the side of the unit, is thus no longer a mere
assumption but follows from the structure; what remains open is only whether the value $74/75$ is
geometrically justified, and $\alpha$ tests this \SM.

The muon lifetime is not an SI pilot either. With $G_F=1/(\sqrt2\,v^2)$,
\begin{equation}
  \frac{\Gamma_\mu}{m_\mu}=\frac{1}{384\,\pi^3}\Bigl(\frac{m_\mu}{v}\Bigr)^4
\end{equation}
up to radiative corrections \QM; $\tau_\mu$ therefore tests only the dimensionless ratio $v/m_\mu$. The
ladder gives $5/(16\,\xi)=2343.75$ for it, the measured value is $2330.34$: a residual of $+43\,\xi$
without $K$ and $+145\,\xi$ with $K$ \KM. Without $K$ the ladder lies $3.4$ times closer, as required by
the rule according to which $K$ has no place in ratios. The residual itself is the known error of the
ladder: relative to $v$, $+82\,\xi$ for the electron, $+43\,\xi$ for the muon and $-34\,\xi$ for the tau,
with the distances of $39$ and $77\,\xi$ between ladder and Koide form \KM. $\tau_\mu$ thus tests the
normalisation of the ladder to $v$, not the factor $K$.

\subsection{The bound hardly needs the Koide amplitude}

The bound in Section~4 assumed both Koide parameters to be exact. The corpus derives the angle
$\theta=2/9$ (Doc.~A110); the justification of the amplitude $d/c=\sqrt2$ is open (R150). If the
amplitude is left free, it absorbs part of the residuals, and only one combination of the mode factors
remains bound. The amplitude, however, shifts almost only the muon-to-electron ratio
($\partial\ln(m_\mu/m_e)/\partial a=32.6$ against $1.8$ for the tau-to-muon ratio) \KM. The bound
combination is therefore almost exactly the difference between tau and muon:

\begin{center}
\begin{tabular}{@{}lll@{}}
\toprule
\textbf{assumed exact} & \textbf{tau--muon} & \textbf{muon--electron}\\
\midrule
$\theta=2/9$ and $d/c=\sqrt2$ & $\le0.92\,\xi$ & $\le0.074\,\xi$\\
only $\theta=2/9$ & $\le0.91\,\xi$ & not bound\\
only $d/c=\sqrt2$ & $\le0.92\,\xi$ & not bound\\
neither & not bound & not bound\\
\bottomrule
\end{tabular}
\end{center}
\KM\ The central statement, that a mode-dependent part of the measurement recursion amounts to at most
about $1\,\%$ of the correction step $100\,\xi$, therefore depends only on the derived angle $2/9$. Only
the finer bound between muon and electron additionally needs the amplitude. If both parameters are
free, there is no bound: two parameters always fit two ratios exactly.

% ============================================================
\section{Consequences for quantum computers}
% ============================================================

\textbf{A common factor is calibrated away.} In the corpus, a gate that is meant to rotate by $\pi$
advances only by $\pi K$, shortened by $100\,\xi=1.33\,\%$ (Docs.~034, 175, 230, 231; A070). R142
records that this deviation is absorbed in the pulse calibration of real qubits and is therefore not
directly measurable. Channel estimation gives the reason: the calibration is a pilot. It measures how
far a pulse actually rotates and adjusts it so that the desired result is obtained. A factor that is
the same for all gates is a single degree of freedom, and the pilot absorbs it completely, just as
$\alpha$ absorbs the SI factor (Section~7). As long as $K$ is common, no quantum computer can see it,
not even with better hardware \BM.

\textbf{Only a mode-dependent part would be visible.} What the calibration cannot absorb is a
difference between modes, for instance between qubits of different frequency or between one- and
two-qubit gates. For the charged leptons such a part is limited to at most about $1\,\xi$
(Sections~4 and 7), that is about $1.2\times10^{-4}$ in relative terms \KM. If a similar bound holds
for qubits, a possible residual signal lies in the same order of magnitude as the fractal damping of
the Bell quantity, $\Delta S/S\approx10^{-5}$ to $10^{-4}$ (R142). The transfer from leptons to qubits
is an analogy and not derived \SM. The gap to the Tsirelson bound $2\sqrt2$ measured today is $0.088$
on ibm\_kingston and $0.104$ pooled (Docs.~147, 384); it is decoherence and about $250$ to $300$ times
larger than a residual signal of this size \KM. This leaves the state of Doc.~384: operationally,
FFGFT and standard quantum mechanics give the same predictions up to a $\xi$ correction that cannot
be resolved with today's hardware.

\textbf{Deterministic reading.} The results support it without proving it. First, the open orbit from
Doc.~394 provides a deterministic mechanism for apparent randomness: a rotation with a non-rational
rotation number never repeats its phases exactly and distributes them uniformly in the long run. With
an advance of $0.00056$ turns per closure cycle, a quarter turn is reached after about $450$ cycles
and a full turn after about $1800$ \KM. This fits the statement of the corpus that noise is
disordered phase (Doc.~384): every individual run is fixed and nevertheless looks like noise over
long runs. The open orbit itself, however, is an interpretation of the residual of $7.6$\,ppm and not
derived \SM. Second, discrete outcomes survive the measurement chain (Section~5): a bit is a winding
number (Doc.~343) and can be regenerated like a digital signal. Third, the Shannon limit remains the
same for both readings: the deterministic individual run and statistical averaging over many runs hit
the same floor (Doc.~343, Theorem~G$'$). The Bell violation is not affected; the FFGFT measurement
remains contextual (Doc.~230).

% ============================================================
\section{Conclusion}
% ============================================================

Communications engineering does not fight the recursion of the channel; it measures it with pilots and
removes it, while the Shannon limit remains. Transferred to FFGFT, this yields a clear picture. In the
ratio sector the channel is flat: with the Koide form as pilot no residual above $1\,\xi$ remains, and
a mode-dependent part of the measurement recursion is limited to about $1\,\%$ of the correction step.
The generation-dependent residuals of the ladder are the error of an inexact pilot, not a channel. In
the SI sector there is exactly one pilot, $\alpha$, and there cannot be a second one: the only relation
with an SI unit fixes the one degree of freedom in which $K$ and the MeV anchoring coincide. The muon
lifetime tests the normalisation of the ladder to $v$, not $K$. Discrete
structures with small denominators survive the measurement chain; this is why the presentation belongs
in the basic form. And a numerical hit counts only if the measurement delivers more information than
the formula contains. For quantum computers this means: a common gate factor disappears in the
calibration, a mode-dependent residual would lie at about $10^{-4}$ and thus far below decoherence;
the deterministic reading remains viable up to the noise limit.

What remains open is whether the value $74/75$, which $\alpha$ tests, is geometrically justified (R135),
and the derivation of the Koide amplitude (R150); only the finer bound between muon and electron still
depends on it \SM.

\medskip
\noindent\textbf{Check script:} \path{2/python/Dok395_Skripte/pruef_395_messkette_kanal.py}

\begin{thebibliography}{99}
\bibitem{Shannon1948} C.~E.~Shannon, \textit{A mathematical theory of communication}, Bell System Technical Journal \textbf{27} (1948) 379--423, 623--656.
\bibitem{PriceGreen1958} R.~Price, P.~E.~Green, \textit{A communication technique for multipath channels}, Proceedings of the IRE \textbf{46} (1958) 555--570.
\bibitem{Brennan1959} D.~G.~Brennan, \textit{Linear diversity combining techniques}, Proceedings of the IRE \textbf{47} (1959) 1075--1102.
\bibitem{WeinsteinEbert1971} S.~B.~Weinstein, P.~M.~Ebert, \textit{Data transmission by frequency-division multiplexing using the discrete Fourier transform}, IEEE Transactions on Communication Technology \textbf{19} (1971) 628--634.
\bibitem{Berrou1993} C.~Berrou, A.~Glavieux, P.~Thitimajshima, \textit{Near Shannon limit error-correcting coding and decoding: Turbo-codes}, Proceedings of ICC '93, Geneva (1993) 1064--1070.
\bibitem{Arikan2009} E.~Arıkan, \textit{Channel polarization}, IEEE Transactions on Information Theory \textbf{55} (2009) 3051--3073.
\bibitem{Koide1982} Y.~Koide, Lettere al Nuovo Cimento \textbf{34} (1982) 201; Physics Letters B \textbf{120} (1983) 161.
\end{thebibliography}
